\chapter{Cancer}

\section*{Patient focus}
Imaging plays a central role in cancer care---from finding a suspicious lump or abnormal blood test, to determining how far cancer has spread, to checking whether treatment is working. Different scans answer different questions: a CT scan shows the size and location of tumors, a PET scan reveals whether cancer is active throughout the body, and an MRI provides detailed images of specific areas like the brain or pelvis. Your doctor will choose the right scan based on the type of cancer and what information is needed.


\label{ch:ch59_cancer}

\section{Role of Imaging}
Imaging supports detection, characterization, staging, treatment planning, response assessment, surveillance, and palliation. The correct study depends on tumor type, suspected site, clinical stage, treatment, and the question being asked.

\subsection{Staging Principles---TNM}
The TNM system is the international standard for cancer staging:

\begin{itemize}
  \item \textbf{T (Primary tumor)}: size, depth of invasion, local extent---e.g.\ T1 small/localized, T4 invasion of adjacent structures
  \item \textbf{N (Regional lymph nodes)}: number, size, and location of involved nodes
  \item \textbf{M (Distant metastasis)}: presence or absence of distant spread
\end{itemize}

Staging is classified as:
\begin{enumerate}
  \item \textbf{Clinical staging (cTNM)}: based on imaging and clinical examination before treatment
  \item \textbf{Pathological staging (pTNM)}: based on surgical specimen---usually more accurate
  \item \textbf{Post-treatment staging (ypTNM)}: after neoadjuvant therapy
\end{enumerate}

\begin{pearlbox}[Imaging determines clinical stage]
For many cancers, the cTNM stage set by imaging determines whether the patient gets surgery, chemotherapy, radiation, or a combination. Accurate staging imaging directly changes treatment pathways. An understaged patient may miss curative surgery; an overstaged patient may receive unnecessary treatment.
\end{pearlbox}

\section{Imaging by Cancer Type}

\begin{table}[h]
\centering
\caption{Imaging strategies for common cancers}
\begin{tabular}{|p{0.14\textwidth}|p{0.30\textwidth}|p{0.20\textwidth}|p{0.20\textwidth}|}
\hline
\textbf{Cancer} & \textbf{Staging} & \textbf{Key sites} & \textbf{PET/CT role} \\
\hline
Breast & Mammo + US, MRI for high risk & Axilla, bone, liver & Limited role \\
Colorectal & CT chest/abdomen/pelvis & Liver, peritoneum & For equivocal CT \\
Liver/HCC & Multiphasic CT/MRI & Vessels, extrahepatic & Limited for HCC \\
Lung & CT chest/abdomen + brain MRI or PET/CT & Brain, adrenal, bone & Excellent for staging \\
Melanoma & CT/PET for stage $>$II & Brain, skin, nodes & Good for high stage \\
Ovarian & CT chest/abdomen/pelvis & Peritoneum, omentum & Limited utility \\
Prostate & MRI pelvis, bone scan & Bone, pelvic nodes & PSMA PET emerging \\
\hline
\end{tabular}
\end{table}

\subsection{Lung Cancer}
\begin{itemize}
  \item \textbf{Detection}: low-dose CT screening (NLST/Lung-RADS criteria)
  \item \textbf{Staging}: PET/CT for whole-body assessment; brain MRI for stage $\geq$II; mediastinal staging with EBUS or mediastinoscopy for PET-positive nodes
  \item \textbf{TNM on imaging}: T by tumor size and invasion (chest wall, mediastinum, great vessels); N by short-axis node diameter ($>$1 cm = suspicious); M by distant lesions
\end{itemize}

\subsection{Breast Cancer}
\begin{itemize}
  \item \textbf{Detection}: mammography (screening and diagnostic), targeted US, breast MRI for high-risk screening and extent-of-disease assessment
  \item \textbf{Staging}: whole-body staging usually with CT chest/abdomen/pelvis; bone scan for bone pain or elevated ALP; MRI brain if neurological symptoms
  \item \textbf{Breast MRI}: useful for multifocal/multicentric disease, contralateral assessment, and neoadjuvant chemotherapy response
\end{itemize}

\subsection{Colorectal Cancer}
\begin{itemize}
  \item \textbf{Staging}: CT chest/abdomen/pelvis with IV contrast; rectal MRI for local staging (circumferential resection margin)
  \item \textbf{Liver metastases}: CT with hepatic arterial phase; MRI with hepatobiliary contrast is more sensitive for small lesions
  \item \textbf{PET/CT}: used when CT is equivocal or to assess operability of isolated metastases
\end{itemize}

\subsection{Prostate Cancer}
\begin{itemize}
  \item \textbf{Staging}: multiparametric MRI (T2, DWI, DCE) for local staging; bone scan for intermediate/high-risk disease; PSMA PET/CT for biochemical recurrence
  \item \textbf{Key MRI features}: extraprostatic extension (bulging capsule, neurovascular bundle involvement), seminal vesicle invasion, tumor volume
\end{itemize}

\section{Response Assessment---RECIST}
\textbf{RECIST 1.1} (Response Evaluation Criteria in Solid Tumors) is the standard framework:

\begin{table}[h]
\centering
\caption{RECIST 1.1 response categories}
\begin{tabular}{|p{0.29\textwidth}|p{0.36\textwidth}|p{0.21\textwidth}|}
\hline
\textbf{Category} & \textbf{Definition} & \textbf{Imaging action} \\
\hline
Complete response (CR) & Disappearance of all target lesions & Confirm at 4 weeks \\
Partial response (PR) & $\geq$30\% decrease in sum of diameters & Continue treatment \\
Progressive disease (PD) & $\geq$20\% increase or new lesion & Reassess treatment \\
Stable disease (SD) & Neither PR nor PD criteria met & Continue treatment \\
\hline
\end{tabular}
\end{table}

\begin{pearlbox}[Pseudoprogression]
Immunotherapy can cause pseudoprogression---apparent tumor growth due to immune cell infiltration rather than true progression. RECIST may be supplemented with immune-related response criteria (irRC, irRECIST) that require confirmed progression on two consecutive scans 4 weeks apart before declaring treatment failure.
\end{pearlbox}

\section{Surveillance After Treatment}
Surveillance imaging aims to detect recurrence early, ideally while still treatable:

\begin{table}[h]
\centering
\caption{Common surveillance intervals (guideline-based, varies by institution)}
\begin{tabular}{|l|l|}
\hline
\textbf{Cancer} & \textbf{Typical surveillance} \\
\hline
Breast & Mammo annually, consider MRI \\
Colorectal & CT chest/abdomen/pelvis annually $\times$5 yr \\
Lung & CT annually $\times$2 yr after treatment \\
Melanoma & Stage-dependent CT/PET \\
Prostate & PSA every 3--6 months, MRI/scan if rising \\
\hline
\end{tabular}
\end{table}

\section{When to Recommend PET/CT}
FDG-PET/CT is most valuable when:
\begin{enumerate}
  \item Staging a cancer known to be FDG-avid (lung, melanoma, lymphoma, oesophageal)
  \item CT is equivocal and biopsy is difficult
  \item Assessing treatment response in lymphoma (Deauville criteria)
  \item Searching for unknown primary with metastatic disease
  \item Suspected recurrence with rising tumor markers but negative conventional imaging
\end{enumerate}

Limitations: false positives (infection, inflammation, granulomatous disease), false negatives (slow-growing tumors, mucinous adenocarcinoma, brain metastases masked by cortical uptake).

\section{Interpretation}
Describe primary lesion, local extension, nodes, distant disease, treatment-related change, and measurable target lesions when applicable. Use validated response criteria only with appropriate technique and timing.

\section{Safety and Communication}
Contrast, radiation, pregnancy, renal function, and biopsy planning require individualized decisions. New cord compression, brain metastasis with mass effect, obstruction, perforation, or major bleeding requires urgent communication.

\begin{pearlbox}[Urgent oncology findings]
Communicate immediately: spinal cord compression (MRI showing cord deformation), superior vena cava obstruction, brain metastases with significant mass effect, pathological fracture risk, airway compromise from tumor, and acute tumor lysis on imaging.
\end{pearlbox}

\clearpage
